\section{The finite-frequency sine bound used at indices 26--30}
\label{ext:finite-sine}
Let the positive and negative support magnitudes lie in
$[a_+,R_+]$ and $[a_-,R_-]$, respectively. Assume
\[
 0<A\le a_+^2+a_-^2,\qquad
 B\ge R_+^2+R_-^2,\qquad M\ge\max(R_+,R_-).
\]
Put $R(u)=1-Au^2(1-Bu^2/42)/20$.
For $u\ge0$, $Mu\le\pi$ and $Bu^2\le42$,
\begin{equation}\label{ext:finite-sine-bound}
 |\Im f(u)|^2\le(\rho^2-\zeta^2)u^6R(u)^2/36.
\end{equation}
In particular $R(u)\ge1-21A/(40B)\ge19/40$.
Outside these frequency gates the ordinary sine-circle inequality
remains available.

Here is a proof retaining the same marginal. Write
$\operatorname{sinc}t=\int_0^1\cos(vt)\,dv$. For $u\sqrt z\le\pi$,
\[
 0\le-\frac{d}{dz}\operatorname{sinc}(u\sqrt z)
 =\frac{u^2}{2}\int_0^1v^2\operatorname{sinc}(uv\sqrt z)\,dv
 \le\frac{u^2}6-\frac{u^4z}{60}+\frac{u^6z^2}{1680}.
\]
The upper bound follows by integrating the global polynomial
$\operatorname{sinc}t\le1-t^2/6+t^4/120$; nonnegativity follows
from $0\le uv\sqrt z\le\pi$.
Integrating between $a^2$ and $b^2$ gives
\[
 |\operatorname{sinc}(au)-\operatorname{sinc}(bu)|
 \le\frac{|a^2-b^2|u^2}{6}
 \left(1-\frac{(a^2+b^2)u^2}{20}
       +\frac{(a^4+a^2b^2+b^4)u^4}{840}\right).
\]
The last bracket is at most $R(u)$, since
$a^4+a^2b^2+b^4\le B(a^2+b^2)$ and
$1-Bu^2/42\ge0$. At $a=b$ use the undivided inequality.

For the centered two-point law at $a,-b$, put
\[
 e_{a,b}=\frac{2a^2b^2}{a+b},\quad
 r_{a,b}=\frac{ab(a^2+b^2)}{a+b},\quad
 h_{a,b}=\frac{abu}{a+b}
       \{\operatorname{sinc}(au)-\operatorname{sinc}(bu)\}.
\]
The identity $(2ab)^2+(a^2-b^2)^2=(a^2+b^2)^2$ gives
$\|(e_{a,b},6h_{a,b}/(u^3R(u)))\|_2\le r_{a,b}$.
Integrate against the exact centered two-point mixture from the
ordinary sine-circle proof. The three integrals are $\zeta$, $\rho$
and $\Im f(u)$, respectively. The triangle inequality proves
\eqref{ext:finite-sine-bound}. The case $u=0$ is immediate.

For each of the $55$ acquired rows ($26\le n\le30$,
$32\le j\le42$), retain its closed probability interval
$[p_l,p_u]$, an actual-noise bound $s\le s_u$, and an upper
bound $d>0$ on the inward sum $a_0+a_1$.
The interface requires $d<p_l$, $d\le1-2p_u$ and
\begin{equation}\label{ext:b1-support-map}
 \frac{p_u^2+(1-p_u-d)^2}{p_u(1-p_u)+s_u}\ge\frac{17}{10},
 \qquad R_-^2+R_+^2\le\frac{1037}{100},
 \qquad \max(R_-,R_+)\le\frac{13}{5}.
\end{equation}
Indeed the minimum of $(p-a_0)^2+(1-p-a_1)^2$ under
$a_0,a_1\ge0$, $a_0+a_1\le d\le1-2p$ is
$p^2+(1-p-d)^2$: use the full inward sum and differentiate in
$a_0$, obtaining $2(1-2p-d)+4a_0\ge0$.
This numerator decreases in $p$ on the retained interval, while
$p(1-p)+s$ increases. Thus \eqref{ext:b1-support-map} gives the
required standardized inner support sum. It licenses
\eqref{ext:finite-sine-bound} with
$A=17/10$, $B=1037/100$, $M=13/5$ on every row.

\paragraph{A direct resource-to-sine map at indices $26$--$30$.}
The three common constants in \eqref{ext:b1-support-map} follow
directly from the acquired floor-$25$ resource.
Write its original simultaneous intercepts as $d_0,L$, and its prices
as $b,g,h$, so that
\[
 d_{\rm floor}-b=g d_0=hL>0,\qquad b\ge0.
\]
These identities are checked exactly against the exported residual;
$L$ is the original joint intercept, even when another bound further
clips the variance. The weighted resource implication gives
\[
 \frac n{25}(b+g\delta)+h\lambda_n\le d_{\rm floor},
 \qquad
 \frac n{25}\frac{\delta}{d_0}+\frac{\lambda_n}{L}\le1.
\]
The second inequality follows by subtracting $nb/25$, dividing by
$d_{\rm floor}-b$, and using $n\ge25$. Therefore
\[
 \delta\le \delta_u:=\min\left\{\frac{25d_0}{n},
                 \frac{25d_{\rm floor}/n-b}{g}\right\}.
\]
The bounded conditional supports give $s\le\mathcal B\delta$, where
$\mathcal B=\max(O_0,O_1)\max(p_u,1-p_l)$ and $\lambda_n=ns$.
Substitution in the joint plane yields
\[
 \lambda_n\le
 \left\{L^{-1}+(25\mathcal B d_0)^{-1}\right\}^{-1}=:\lambda_u.
\]
This step uses one joint plane before optimizing. For all $55$ rows,
the choices $d=\delta_u$ and $s_u=\lambda_u/n$ satisfy the two
preceding gates and
\[
 \frac{p_u^2+(1-p_u-\delta_u)^2}
      {p_u(1-p_u)+\lambda_u/n}>\frac{17}{10}.
\]
The smallest left side is greater than $1716/1000$ and occurs
at $n=26$ on $P_{38}$. The same rows satisfy
$R_-^2+R_+^2\le1037/100$ and $\max(R_-,R_+)\le13/5$.
The actual-index support proofs and the full $100$-panel resource
cover are checked separately. Thus the finite-sine majorant applies
on every surviving panel. Its integral uses only these common
constants and the rational moment and frequency intervals below.

The symmetrization inputs admit a shorter derivation than the older
special first-moment envelope. Since
$(13/5)^2< (11/10)^2+8$, the general bounded-support inequality
for $\eta$ gives $\eta\le b(\rho)$ throughout $[11/10,6/5]$,
where $b(r)=(\sqrt{r^2+8}-r)/2$.
The universal inequality $\zeta\ge b(\rho)$ was proved separately.
Since $b$ decreases and $r+b(r)$ increases, on a strip $[r_l,r_u]$
we may use simultaneously
\[
 \rho+\eta\le r_u+b(r_u),\qquad
 \rho^2-\zeta^2\le r_u^2-b(r_u)^2.
\]
Rational brackets $z_l,z_u$ for $b(r_u)$ are checked by
$z_l^2+r_uz_l\le2\le z_u^2+r_uz_u$ with $0<z_l\le z_u$.
The integrator uses $r_u+z_u$ and $r_u^2-z_l^2$ in their indicated
upper-bound directions. The value $z_u$ is not separately asserted
to bound $\eta$ at the lower end of the strip.

The selected $23$ closed strips cover $[11/10,6/5]$. All five
actual indices are evaluated on each strip, giving $115$ complete
four-term Fourier inequalities. The optional fourth-Taylor real
endpoint and Lemma~\ref{ext:two-corners} are used only when their
whole-cell gates hold. The remaining cells use a previously valid
geometric-sum or direct interval bound.

The implementation sometimes integrates the maximum of the two
corner majorants more accurately. If $f_0,f_1$ are twice continuously
differentiable on $[t_0-h,t_0+h]$, enclose their midpoint values above
by $a_i$, their midpoint derivatives by $[d_i^-,d_i^+]$, and both
second derivatives above by $C\ge0$. Taylor's theorem bounds the
integral of their pointwise maximum by
\[
 \int_0^h\max_i(a_i+d_i^+x)\,dx
 +\int_0^h\max_i(a_i-d_i^-x)\,dx+\frac{Ch^3}{3}.
\]
Each affine maximum has at most one crossing, found exactly. This
bounds the maximum before integration and assumes nothing about
how often the original smooth curves cross.
